\section{The split space and its Clifford algebra}
\label{sec:arena}

\subsection{A fixed realisation}

Let $V=\R^{3,3}$ with orthonormal basis $e_1,\dots,e_6$, $e_a^2=\eta_{aa}$ and $\eta=\operatorname{diag}(1,1,1,-1,-1,-1)$.
We call $V^+=\langle e_1,e_2,e_3\rangle$ the temporal triple and $V^-=\langle e_4,e_5,e_6\rangle$ the spatial triple. The
Clifford algebra $\Cl(3,3)$ is generated by the $e_a$ subject to $e_ae_b+e_be_a=2\eta_{ab}\id$; it is isomorphic to
$M_8(\R)$, and we fix once and for all its realisation on $S=\R^8=(\R^2)^{\otimes3}$,
\begin{equation}
\label{eq:realisation}
\begin{aligned}
e_1&=\sigma_x\otimes\id\otimes\id, & e_2&=\sigma_z\otimes\sigma_x\otimes\id, & e_3&=\sigma_z\otimes\sigma_z\otimes\sigma_x,\\
e_4&=\varepsilon\otimes\id\otimes\id, & e_5&=\sigma_z\otimes\varepsilon\otimes\id, & e_6&=\sigma_z\otimes\sigma_z\otimes\varepsilon,
\end{aligned}
\end{equation}
with $\sigma_x=\left(\begin{smallmatrix}0&1\\1&0\end{smallmatrix}\right)$,
$\sigma_z=\left(\begin{smallmatrix}1&0\\0&-1\end{smallmatrix}\right)$ and
$\varepsilon=\left(\begin{smallmatrix}0&-1\\1&0\end{smallmatrix}\right)$. All six matrices are signed permutation
matrices. For a subset $I=\{a_1<\dots<a_r\}$ we write $e_I=e_{a_1}\cdots e_{a_r}$ for the corresponding blade of grade $r$;
the $64$ blades form a basis of $M_8(\R)$.

The tower is built by graded tensor products. At $n$ constituents the spinor module is $S_n=(\R^8)^{\otimes n}$ and the
generators of the $k$-th constituent are
\begin{equation}
\label{eq:tower}
e^{(k)}_a=\underbrace{\omega\otimes\cdots\otimes\omega}_{k-1}\otimes e_a\otimes\id\otimes\cdots\otimes\id,
\qquad \omega:=e_1e_2\cdots e_6 .
\end{equation}
Since $\omega^2=\id$ and $\omega$ anticommutes with every $e_a$, generators of different constituents anticommute, and the
$6n$ matrices $e^{(k)}_a$ generate $\Cl(3n,3n)\cong M_{8^n}(\R)$. When we need a single index we write
$E_{6(k-1)+a}=e^{(k)}_a$, so that at three constituents the eighteen generators are $E_1,\dots,E_{18}$ and the temporal
ones are $E_A$ with $A\in\{1,2,3,7,8,9,13,14,15\}$.

Two elementary rules carry most of the finite-dimensional arguments of this paper.

\begin{lemma}[Blade rules]
\label{lem:blade}
For blades $e_I,e_J$ of grades $|I|,|J|$,
\begin{equation}
\label{eq:bladecomm}
e_Ie_J=(-1)^{|I||J|-|I\cap J|}\,e_Je_I ,
\end{equation}
and for a blade $B$ of grade $r$
\begin{equation}
\label{eq:transpose}
B^{\!\top}=(-1)^{r(r-1)/2}\Bigl(\prod_{a\in I}\eta_{aa}\Bigr)B=B^2B .
\end{equation}
In particular a blade is antisymmetric if and only if it squares to $-\id$, and symmetric if and only if it squares to $+\id$.
\end{lemma}

\begin{proof}
For \eqref{eq:bladecomm}, move each generator of $e_J$ past the generators of $e_I$: it anticommutes with the $|I|$ of them
it differs from and commutes with itself when it occurs in $I$, which gives the exponent. For \eqref{eq:transpose}, the
matrices $\sigma_x,\sigma_z$ are symmetric and $\varepsilon$ is antisymmetric, so each $e_a$ satisfies $e_a^{\!\top}=\eta_{aa}e_a$;
transposing a product reverses it, and restoring the order costs $(-1)^{r(r-1)/2}$. The same two factors give
$B^2=(-1)^{r(r-1)/2}\prod_{a\in I}\eta_{aa}\,\id$, so $B^{\!\top}=B^2B$.
\end{proof}

\begin{example}
\label{ex:blades}
The bivector $e_2e_3$ has $(e_2e_3)^2=-e_2^2e_3^2=-\id$, so it is antisymmetric and $\exp(\theta e_2e_3)=\cos\theta\,\id+
\sin\theta\,e_2e_3$ is a rotation of period $2\pi$. The bivector $e_1e_4$ has $(e_1e_4)^2=-e_1^2e_4^2=+\id$, so it is
symmetric and $\exp(\theta e_1e_4)=\cosh\theta\,\id+\sinh\theta\,e_1e_4$ is a boost, unbounded in $\theta$. The two
bivectors share the index $1$ and $|I||J|-|I\cap J|=4-1=3$, so they anticommute; $e_2e_3$ and $e_5e_6$ share nothing and
commute.
\end{example}

\subsection{The signature rule}
\label{sec:sigrule}

\begin{theorem}[Gauging versus grading]
\label{thm:sig}
For $a\neq b$, $(e_ae_b)^2=-e_a^2e_b^2$. A bivector therefore squares to $-\id$ when its two directions have the same
sign and to $+\id$ when they have opposite signs, and by \cref{lem:blade} the first kind are exactly the antisymmetric
bivectors. The three time--time and three space--space bivectors span the maximal compact subalgebra
$\so3\oplus\so3\subset\so{3,3}$, and the nine mixed bivectors are boosts.
\end{theorem}

\begin{corollary}[No gauge boson is a boost]
\label{cor:noboost}
For a mixed bivector $B$, $\exp(\theta B)=\cosh\theta\,\id+\sinh\theta\,B$ generates $\R$ and not $U(1)$: it carries no
phase and defines no compact charge.
\end{corollary}

The criterion that matters for a gauge generator is antisymmetry with respect to the positive form that a kinetic term
supplies, equivalently membership of the $+1$ eigenspace of the Cartan involution $\theta(X)=-X^{\!\top}$
\cite{Knapp,Helgason}. For a single blade this coincides with $X^2=-\id$ by \eqref{eq:transpose}, which is why the rule is
decisive for bivectors. For general elements it is not: $\operatorname{diag}(J_2,2J_2)$, with $J_2$ the standard $2\times2$
rotation generator, generates a group of period $2\pi$ although its square is $\operatorname{diag}(-\id,-4\id)$.

\begin{figure}[t]
\centering
\begin{tikzpicture}[scale=0.82]
\foreach \i in {1,...,6}{
  \node at (\i-0.5,0.35) {\small $e_{\i}$};
  \node at (-0.35,-\i+0.5) {\small $e_{\i}$};
}
\foreach \i in {1,...,6}{
  \foreach \j in {1,...,6}{
    \pgfmathtruncatemacro{\ti}{\i<4 ? 1 : 0}
    \pgfmathtruncatemacro{\tj}{\j<4 ? 1 : 0}
    \ifnum\j>\i
      \ifnum\ti=\tj
        \fill[tcol!35] (\j-1,-\i+1) rectangle (\j,-\i);
      \else
        \fill[scol!25] (\j-1,-\i+1) rectangle (\j,-\i);
      \fi
    \fi
    \ifnum\j=\i \fill[lgray] (\j-1,-\i+1) rectangle (\j,-\i); \fi
  }
}
\draw[step=1,mgray,thin] (0,-6) grid (6,0);
\draw[very thick] (0,0) rectangle (6,-6);
\draw[thick] (3,0) -- (3,-6); \draw[thick] (0,-3) -- (6,-3);
% pole operators
\foreach \i in {1,2,3}{\pgfmathtruncatemacro{\jj}{\i+3}\node at (\jj-0.5,-\i+0.5) {\footnotesize $\omega_{\i}$};}
\node at (2.5,-1.5) {\footnotesize $J$};
\fill[tcol!35] (6.6,-1.05) rectangle ++(0.32,0.32); \draw[mgray] (6.6,-1.05) rectangle ++(0.32,0.32);
\node[align=left,anchor=west] at (7.05,-0.9) {\small same sign: $B^2=-\id$, $B^{\!\top}=-B$,};
\node[align=left,anchor=west] at (7.05,-1.5) {\small compact; $3+3$ generators of $\so3\oplus\so3$};
\fill[scol!25] (6.6,-2.65) rectangle ++(0.32,0.32); \draw[mgray] (6.6,-2.65) rectangle ++(0.32,0.32);
\node[align=left,anchor=west] at (7.05,-2.5) {\small mixed sign: $B^2=+\id$, $B^{\!\top}=B$,};
\node[align=left,anchor=west] at (7.05,-3.1) {\small boosts; $9$ generators};
\node[align=left,anchor=west] at (6.6,-4.1) {\small $\omega_i=e_ie_{i+3}$: the pole operators,};
\node[align=left,anchor=west] at (6.6,-4.7) {\small \phantom{$\omega_i$} commuting involutions};
\node[align=left,anchor=west] at (6.6,-5.5) {\small $J=e_2e_3$: the interior charge};
\end{tikzpicture}
\caption{The fifteen bivectors $e_ae_b$, $a<b$, of $\Cl(3,3)$. Blue cells pair two directions of the same sign and generate
compact rotations; red cells pair a temporal with a spatial direction and generate boosts (\cref{thm:sig}). The diagonal of
the red block holds the three pole operators, which grade the algebra (\cref{prop:grading}); the cell $(2,3)$ holds the
bivector $J$ that will become the electromagnetic charge and the imaginary unit.}
\label{fig:bivectors}
\end{figure}

\subsection{The pole operators and a multigrading}
\label{sec:grading}

The mixed bivectors are not idle. The three \emph{pole operators}
\[
\omega_i:=e_ie_{i+3}\qquad(i=1,2,3)
\]
are boosts, so $\omega_i^2=+\id$, and they commute pairwise by \eqref{eq:bladecomm}, since any two share no index. Square
$-\id$ gauges and square $+\id$ grades, and no bivector does both.

\begin{proposition}[A $\Z_2^3$ grading of $\Cl(3,3)$]
\label{prop:grading}
For a blade $e_I$ put $\chi_i(e_I)=(-1)^{|I\cap\{i,i+3\}|}$. Then $\omega_ie_I\omega_i^{-1}=\chi_i(e_I)\,e_I$, and
$\Cl(3,3)=\bigoplus_{\sigma\in\{\pm\}^3}\Cl_\sigma$ is a multiplicative grading with eight sectors of dimension eight.
\begin{enumerate}[label=(\roman*),leftmargin=2.2em]
\item The joint spectrum of $(\omega_1,\omega_2,\omega_3)$ on $S$ consists of all eight sign triples, each on a line, so the
pole operators cut $\R^8$ into eight one-dimensional eigenlines.
\item The vectors occupy the three weights with exactly one minus sign, two each ($e_i$ and $e_{i+3}$); the bivectors occupy
the four even weights with multiplicities $3,4,4,4$; the all-minus weight is spanned by eight trivectors, namely $e_1e_2e_3$,
$e_4e_5e_6$ and six mixed ones.
\item Correspondingly $\so{3,3}=\mathfrak g_0\oplus\mathfrak g_1\oplus\mathfrak g_2\oplus\mathfrak g_3$ is graded by the Klein
four-group, of dimensions $3+4+4+4$, with $\mathfrak g_0=\operatorname{span}\{\omega_1,\omega_2,\omega_3\}$ abelian,
self-normalising and consisting of semisimple elements with real eigenvalues: a maximally split Cartan subalgebra of
$\so{3,3}\cong\slr{4,\R}$.
\end{enumerate}
\end{proposition}

\begin{proof}
The conjugation sign is \eqref{eq:bladecomm} with $|J|=2$: $(-1)^{2|I|-|I\cap\{i,i+3\}|}$. The weight is multiplicative because
$|I\triangle I'|$ and $|I|+|I'|$ have the same parity in each pair. The census in (ii) is a count of how many generators of
each pair a blade contains. For (iii), a bivector inside one pair is a pole operator, of weight $(+,+,+)$; a bivector with one
generator in each of two pairs has two minus signs, and each of the three such weights holds $2\times2=4$ bivectors. The
pole operators commute, are symmetric with eigenvalues $\pm1$, and the centraliser of $\mathfrak g_0$ in $\so{3,3}$ is
$\mathfrak g_0$ itself, because every other bivector has a nontrivial weight; its dimension equals the real rank three of
$\so{3,3}$.
\end{proof}

The three characters are grading weights, not metric signatures: $e_1$ and $e_4$ carry the same weight and opposite squares,
so nothing in the grading tells a temporal from a spatial direction. The decomposition in (iii) is the root decomposition of
$\slr{4,\R}$ with respect to a split Cartan subalgebra, grouped by the root lattice modulo two; gradings of Clifford and
classical Lie algebras by $\Z_2^n$ are well known \cite{MosnaMirallesVaz}, and we use this one because the pole operators will
reappear in the field-strength reading as the six planes $(e_i,e_{i+3})$ that carry one complex component each.

\begin{example}[The paired-mode projector]
\label{ex:pairedmode}
Let $\mathsf P=\omega_1+\omega_2+\omega_3$ be the sum of the pole operators and $\mathcal V=e_1e_2\cdots e_6$. On the joint
eigenline of signs $(\varepsilon_1,\varepsilon_2,\varepsilon_3)$, $\mathsf P=\varepsilon_1+\varepsilon_2+\varepsilon_3\in\{3,1,-1,-3\}$
with multiplicities $1,3,3,1$. Reordering $\omega_1\omega_2\omega_3=e_1e_4e_2e_5e_3e_6$ into $e_1\cdots e_6$ costs three
transpositions, so $\mathcal V=-\varepsilon_1\varepsilon_2\varepsilon_3$ on that line. Hence $\tfrac12(\mathsf P+\mathcal V)$ takes the
values $1,1,-1,-1$ on the four strata: it is the sign of $\mathsf P$, and equals the polynomial $(13\mathsf P-\mathsf P^3)/12$.
The operator
\begin{equation}
\label{eq:Piclosed}
\Pi:=\tfrac12\id-\tfrac14(\mathsf P+\mathcal V)
\end{equation}
is therefore the orthogonal projector of rank four onto the strata with $\mathsf P<0$. We call it the paired-mode projector; it
is a natural candidate for an observation projector at one constituent and we test it against the photon in
\cref{sec:photon}.
\end{example}

\subsection{The interior four-plane and the surplus plane}
\label{sec:interior}

The physical reading of the arena singles out a Lorentzian four-plane inside $V$. We take
\[
E=\langle e_1,e_4,e_5,e_6\rangle,\qquad \gamma^0=e_1,\quad \gamma^a=e_{a+3}\ (a=1,2,3),
\]
of signature $(1,3)$, and call the orthogonal complement $U=\langle e_2,e_3\rangle$ the surplus plane. The interior Lorentz
algebra $\lint$ is spanned by the six bivectors of $E$, with rotations $S_1=\tfrac12\gamma^2\gamma^3$ (cyclically) and boosts
$T_a=\tfrac12\gamma^0\gamma^a$. Its spin Casimir on $S$ is $-\sum_aS_a^2=\tfrac34\id$, and a rotation through $2\pi$ is $-\id$
on $S$ while it is $+\id$ on vectors (\cref{sec:spin}): the same rotors give spin one-half to spinors and spin one to vectors.

\begin{proposition}[The interior charge is unique]
\label{prop:6plus8plus1}
As $\lint$-modules, $\Lambda^2V=\Lambda^2E\oplus(E\otimes U)\oplus\Lambda^2U$, of dimensions $6+8+1$, is the adjoint
representation, two copies of the vector representation, and the trivial representation. The centraliser of $\lint$ in
$\so{3,3}$ is therefore the line spanned by
\[
J:=e_2e_3 ,
\]
which satisfies $J^2=-\id$, $J^{\!\top}=-J$, and commutes with each of the four gamma matrices.
\end{proposition}

\begin{proof}
The adjoint and vector representations of $\so{1,3}$ contain no invariant, so only $\Lambda^2U$ survives. That $J$ commutes
with $e_1,e_4,e_5,e_6$ is \eqref{eq:bladecomm} with disjoint index sets and $|I||J|=2$.
\end{proof}

\begin{figure}[t]
\centering
\begin{tikzpicture}[>=Stealth]
\coordinate (t1) at (0,2.6); \coordinate (t2) at (0,1.3); \coordinate (t3) at (0,0);
\coordinate (s1) at (5,2.6); \coordinate (s2) at (5,1.3); \coordinate (s3) at (5,0);
% interior four-plane (L-shaped region)
\fill[gcol!8,rounded corners=8pt] (-1.85,3.05) -- (7.0,3.05) -- (7.0,-0.5) -- (4.45,-0.5) -- (4.45,2.0) -- (-1.85,2.0) -- cycle;
\draw[gcol,very thick,rounded corners=8pt] (-1.85,3.05) -- (7.0,3.05) -- (7.0,-0.5) -- (4.45,-0.5) -- (4.45,2.0) -- (-1.85,2.0) -- cycle;
% surplus plane
\draw[ocol,very thick,dashed,rounded corners=6pt] (-0.75,-0.5) rectangle (0.55,1.75);
% pairing arrows
\draw[->,gray,thick,shorten <=4pt,shorten >=4pt] (t1) -- node[above] {\small $e_i\mapsto e_{i+3}$} (s1);
\draw[->,gray,thick,shorten <=4pt,shorten >=4pt] (t2) -- (s2);
\draw[->,gray,thick,shorten <=4pt,shorten >=4pt] (t3) -- (s3);
\foreach \p in {t1,t2,t3}{\fill[tcol] (\p) circle (2.4pt);}
\foreach \p in {s1,s2,s3}{\fill[scol] (\p) circle (2.4pt);}
\node[left=3pt] at (t1) {$e_1=\gamma^0$}; \node[left=3pt] at (t2) {$e_2$}; \node[left=3pt] at (t3) {$e_3$};
\node[right=3pt] at (s1) {$e_4=\gamma^1$}; \node[right=3pt] at (s2) {$e_5=\gamma^2$}; \node[right=3pt] at (s3) {$e_6=\gamma^3$};
\node[tcol,anchor=south] at (0,3.15) {\small $V^+$ (times): $e_i^2=+\id$};
\node[scol,anchor=south] at (5,3.15) {\small $V^-$ (lengths): $e_{i+3}^2=-\id$};
\node[ocol,anchor=east,align=right] at (-0.95,0.6) {\small surplus plane $U$,\\ \small $J=e_2e_3$};
\node[gcol,anchor=north,align=center] at (5.7,-0.6) {\small interior four-plane\\ \small $E=\langle e_1,e_4,e_5,e_6\rangle$, signature $(1,3)$};
\end{tikzpicture}
\caption{The arena. The positive triple $V^+$ (times) and the negative triple $V^-$ (lengths) are paired by $e_i\mapsto e_{i+3}$;
the pairing has two square roots, the grading $\cK$ and the binding operator $G$ (\cref{prop:twoK}). The interior four-plane
$E$, on which the Dirac sector is built, contains the clock $e_1$ and the whole spatial triple; its complement is the surplus
plane $U=\langle e_2,e_3\rangle$, whose rotation generator $J$ is the only bivector commuting with the interior Lorentz algebra
(\cref{prop:6plus8plus1}).}
\label{fig:arena}
\end{figure}

\Cref{prop:6plus8plus1} is the first instance of a pattern that recurs throughout: a structure is fixed once the right
ambient algebra and the right conditions are named, and it ceases to be fixed when either is relaxed. Here the ambient
space is the span of the bivectors and the condition is commutation with $\lint$. In \cref{sec:dirac} we enlarge the ambient
algebra to all of $\End_\R(S)$ and find that a different condition, commutation with the internal bivectors, is what keeps the
answer unique.
